% APPROACH HINT: Apply the Lonely Points Lemma: for each s in S, choose a linear functional uniquely maximized at s (this exists since s is a vertex of conv(S)), then maximize it over s+T to obtain a point ell_s = s + t_s that cannot come from any other s' in S. This gives a unique representation (claim 1). The points ell_s are distinct because a single point cannot simultaneously maximise two different linear functionals that separate distinct convex vertices (claim 2). Since ell_s has only one representation, F(ell_s) = epsilon(t_s) which is non-zero (claim 3). This immediately yields |{p : F(p) != 0}| >= n (claim 4). For the equality case (claim 5): if exactly n points have F != 0 and the n lonely points all have F != 0, then the support equals {ell_s : s in S} and every other point in S+T must have F = 0, i.e., all its representations cancel. The output must be a self-contained LaTeX fragment (no \documentclass header), just theorem/proof environments with explicit \begin{proof}...\end{proof} blocks.
\documentclass{article}
\usepackage{amsmath,amssymb,amsthm}
\begin{document}

\textbf{Subproblem 1 --- Lonely-points lower bound.}

\medskip
\textbf{Statement.}
Let $S\subset\mathbb{R}^2$ be a finite set of $n$ points in convex position (every point is a vertex of $\operatorname{conv}(S)$).
Let $T\subset\mathbb{R}^2$ be a finite set of translation vectors and let $\varepsilon:T\to\{+1,-1\}$ be a sign function.
Define
\[
  F(p) \;=\; \sum_{\substack{t\in T\\ p-t\in S}} \varepsilon(t),\qquad p\in\mathbb{R}^2.
\]
Prove:
\begin{enumerate}
  \item For every $s\in S$ there exists $t_s\in T$ such that the point $\ell_s := s+t_s$ has a \emph{unique} representation in $S+T$, i.e.\ $\ell_s - t = s$ for every $t\in T$ with $\ell_s - t \in S$ implies $t = t_s$.  (Such a point is called \emph{lonely}.)
  \item The $n$ lonely points $\ell_{s_1},\ldots,\ell_{s_n}$ are pairwise distinct.
  \item $F(\ell_s) = \varepsilon(t_s)\neq 0$ for every $s\in S$.
  \item Consequently $|\{p\in\mathbb{R}^2 : F(p)\neq 0\}| \ge n$.
  \item If moreover $|\{p : F(p)\neq 0\}| = n$, then $\{p : F(p)\neq 0\} = \{\ell_s : s\in S\}$ and every point in $S+T$ that is \emph{not} of the form $\ell_s$ satisfies $F(p)=0$.
\end{enumerate}

\end{document}
